\documentclass[12pt]{article}
\usepackage[utf8]{inputenc}
\usepackage[T1]{fontenc}
\usepackage{hyperref}
\title{Meänkieli Language Resources for K-12 Curriculum Development}
\author{Compiled by Claude Sonnet 5.5 (claude-sonnet-5-5) for Andr\'as Kornai}
\date{}
\begin{document}
\maketitle

\section{Introduction}
Meänkieli (ISO 639-3 fit) is a Finnic language of the Torne Valley in Sweden, a recognised national minority language there. Its standard written form is based on the dialects around Pajala, Övertorneå and Haparanda. Note that the reference description supplied with the task (older adults only, not known to be taught in schools) is at odds with the search results below, which describe school curricula and textbooks in Sweden. This report is based on a small number of web searches and is not exhaustive.

\section{Primary Educational Resources}
\begin{itemize}
\item Sveriges Radio report (2015) on Skolverket's plan to give Meänkieli and the other minority languages their own curricula and teaching materials, noting a teacher shortage: \url{https://www.sverigesradio.se/artikel/6133178}
\item Umeå University syllabus page for a Meänkieli course; the search listing associated the Umeå course literature with the grade 6--9 textbook \emph{Tule sieki} and with \emph{Meänkieli} by Nylund, Nylund and Nylund (Pajala kommun, 2007). I did not open the textbooks themselves: \url{https://www.umu.se/en/student/syllabus/1fi060/}
\item ECML document discussing minority-language education, including the first introduction of Meänkieli as a school subject in 1999: \url{https://www.ecml.at/mtp2/Ensemble/pdf/A2Webtext_Huss_F.pdf}
\end{itemize}

\section{Linguistic and Reference Materials}
\begin{itemize}
\item Meänkieli--Swedish--Meänkieli dictionary (Elina Kangas, Swedish Language Council, with the National Association of Swedish Tornedalians), about 33,000 words, data released as XML under CC0: \url{https://snd.gu.se/en/catalogue/study/2022-241}
\item Wikipedia, Meänkieli: \url{https://en.wikipedia.org/wiki/Me%C3%A4nkieli}
\item Wikipedia, Meänkieli grammar: \url{https://en.wikipedia.org/wiki/Me%C3%A4nkieli_grammar}
\item Blokland and Valijärvi, ``Aspects of Meänkieli from a grammaticographical perspective'': \url{https://www.keeljakirjandus.ee/en?p=31859}
\end{itemize}

\section{Religious and Community Resources}
No religious or community-organisation resources were found in my searches, other than the Tornedalians' association named as co-compiler of the dictionary above.

\section{Digital and Audio Resources}
The dictionary data (above) is machine-readable. Valijärvi's article on language ideologies in Meänkieli teaching and learning resources (UCL Discovery) discusses digital and interactive materials: \url{https://discovery.ucl.ac.uk/id/eprint/10215021}. No audio or video collections were found in these searches.

\section{Curriculum Development Context}
Sweden has moved toward national curricula and subject materials for Meänkieli; the main reported obstacle is a shortage of qualified teachers. A usable base for a curriculum exists (textbooks, dictionary, descriptive literature), but I could not verify the current availability or content of the textbooks.

\end{document}